\section{Gaussian filtering between two Hamiltonians}
\label{sec:gaussian_filter}

Let $\mathcal H$ be a finite-dimensional complex Hilbert space, with inner
product conjugate-linear in its first argument. All operator norms below
are induced by the Hilbert-space norm. The Gaussian-filter argument in the
proof of the reset lemma of~\cite{OpenAI2026PolynomialPEPS} averages
$e^{it\mathbb H'}We^{-it\mathbb H}$ over time. Its variance $h$ determines
both the suppression of excited energies and the cost of truncating time.
% Source: polynomial-PEPS manuscript (September 24, 2026), lem:reset,
% pinned peps-02-information-adc7f124.tex, lines 426--452; the real overlap
% is established separately at eq:info-reset-overlap, lines 402--414.
% This leaf covers the finite-dimensional integral and gap argument only.
% It makes no spatial-locality, full-reset, or Hermiticity claim for M.
% The source passage is algebraic; no tensor diagram is needed.

\subsection{The normalized Gaussian and its tail}

\begin{definition}[Gaussian time density]
    \label{def:gaussian_time_density}
    \lean{GaussianFilter.kernel, GaussianFilter.kernel_eq,
        GaussianFilter.kernel_zero}
    \leanok
    For $h>0$, define the centered Gaussian density of variance $h$ by
    \begin{align}
        g_h(t)&=(2\pi h)^{-1/2}\exp\!\left(-\frac{t^2}{2h}\right).
        \label{eq:gaussian_density}
    \end{align}
    Write $\mu_h$ for the corresponding probability measure. At $h=0$,
    set $\mu_0=\delta_0$ and $g_0(t)=0$ for every $t$.
    Thus $g_0$ is not a density of $\mu_0$ with respect to Lebesgue measure.
\end{definition}

\begin{theorem}[Positivity and normalization]
    \label{thm:gaussian_normalization}
    \lean{GaussianFilter.kernel_nonneg, GaussianFilter.integrable_kernel,
        GaussianFilter.integral_kernel}
    \leanok
    \uses{def:gaussian_time_density}
    For every $h\ge0$, the function $g_h$ is nonnegative and integrable.
    For $h>0$, $\int_{\R}g_h(t)\,dt=1$.
\end{theorem}
\begin{proof}\leanok
    \uses{def:gaussian_time_density}
    For positive variance these are the positivity, integrability, and
    normalization identities for the centered Gaussian density. At zero
    variance, $g_0=0$.
\end{proof}

\begin{theorem}[Density and probability integrals]
    \label{thm:gaussian_density_integrals}
    \lean{GaussianFilter.integrable_kernel_smul_iff,
        GaussianFilter.integral_kernel_smul,
        GaussianFilter.setIntegral_kernel_smul,
        GaussianFilter.integral_kernel_eq_measureReal}
    \leanok
    \uses{def:gaussian_time_density}
    Let $h>0$ and let $f$ take values in a real normed vector space.
    Then $g_hf$ is Lebesgue integrable if and only if $f$ is integrable
    with respect to $\mu_h$. The integrals satisfy
    \begin{align}
        \int_{\R}g_h(t)f(t)\,dt&=\int_{\R}f(t)\,d\mu_h(t),
        \label{eq:gaussian_density_integral}\\
        \int_S g_h(t)f(t)\,dt&=\int_S f(t)\,d\mu_h(t) \notag
    \end{align}
    for every measurable set $S$. In particular,
    $\int_S g_h(t)\,dt=\mu_h(S)$.
\end{theorem}
\begin{proof}\leanok
    \uses{def:gaussian_time_density}
    Substitute $d\mu_h(t)=g_h(t)\,dt$ in the integrability criterion and
    the integral. Apply the same identity to $\mathbf 1_Sf$ for the
    restricted integral. The scalar mass identity is the density formula
    for $\mu_h$.
\end{proof}

\begin{theorem}[Characteristic integral and exponential moments]
    \label{thm:gaussian_characteristic}
    \lean{GaussianFilter.integral_cexp,
        GaussianFilter.integral_kernel_cexp,
        GaussianFilter.hasSubgaussianMGF_gaussianReal}
    \leanok
    \uses{def:gaussian_time_density}
    For $h\ge0$ and $\omega,s\in\R$,
    \begin{align}
        \int_{\R}e^{i\omega t}\,d\mu_h(t)&=e^{-h\omega^2/2},
        \label{eq:gaussian_characteristic}\\
        \int_{\R}e^{st}\,d\mu_h(t)&\le e^{hs^2/2}. \notag
    \end{align}
    In particular, all these exponential moments are finite, and the
    centered Gaussian is sub-Gaussian with variance parameter $h$.
    For $h>0$, the first integral also equals
    $\int_{\R}g_h(t)e^{i\omega t}\,dt$.
    The frequency $\omega$ is angular frequency, so no additional factor
    of $2\pi$ occurs in the exponent.
\end{theorem}
\begin{proof}\leanok
    \uses{thm:gaussian_density_integrals}
    The characteristic function gives the first identity. The exact
    Gaussian moment-generating function is $e^{hs^2/2}$, which implies
    the exponential-moment and sub-Gaussian bounds. For $h>0$, use
    (\ref{eq:gaussian_density_integral}) to obtain the density integral.
\end{proof}

For $T\ge0$, write $I_T=[-T,T]$. Set
$\tau_h(T)=2\exp(-T^2/(2h))$ for $h>0$, and $\tau_0(T)=2$.
The value at zero permits a common bound for all nonnegative variances;
the actual mass of $I_T^c$ under $\mu_0=\delta_0$ is zero.

\begin{theorem}[Discarded Gaussian mass]
    \label{thm:gaussian_tail}
    \lean{GaussianFilter.measure_compl_Icc_le,
        GaussianFilter.integral_kernel_compl_Icc_le}
    \leanok
    \uses{def:gaussian_time_density}
    For $h\ge0$ and $T\ge0$, $\mu_h(I_T^c)\le\tau_h(T)$.
    For $h>0$, this gives the density estimate
    \begin{align}
        \int_{I_T^c}g_h(t)\,dt&\le 2e^{-T^2/(2h)}.
        \label{eq:gaussian_density_tail}
    \end{align}
\end{theorem}
\begin{proof}\leanok
    \uses{thm:gaussian_characteristic, thm:gaussian_density_integrals}
    The exponential-moment bound applied to $t$ and $-t$ gives
    $\mu_h\{t\ge T\}\le e^{-T^2/(2h)}$ and
    $\mu_h\{-t\ge T\}\le e^{-T^2/(2h)}$ when $h>0$.
    Their sum bounds $\mu_h(I_T^c)$ by the union bound. At $h=0$,
    the discarded mass is zero. The density identity follows from
    Theorem~\ref{thm:gaussian_density_integrals}.
\end{proof}

\subsection{The two-Hamiltonian integral}

\begin{definition}[Full and truncated Gaussian filters]
    \label{def:gaussian_filter_integrals}
    \lean{GaussianFilter.gaussianIntertwiner,
        GaussianFilter.gaussianIntertwinerTruncated}
    \leanok
    \uses{def:gaussian_time_density}
    For matrices $\mathbb H',\mathbb H,W$ on $\mathcal H$, define
    $A_t=e^{it\mathbb H'}We^{-it\mathbb H}$ and, for $h\ge0$,
    \begin{align}
        M_h(\mathbb H',\mathbb H;W)&=\int_{\R}A_t\,d\mu_h(t),
        \label{eq:gaussian_filters}\\
        M_h^T(\mathbb H',\mathbb H;W)&=\int_{[-T,T]}A_t\,d\mu_h(t). \notag
    \end{align}
    Truncation retains the original Gaussian mass, without division
    by $\mu_h([-T,T])$.
    Write $M_h$ and $M_h^T$ when the three matrices are fixed.
\end{definition}

\begin{theorem}[Continuity, integrability, and constant norm]
    \label{thm:gaussian_integrand}
    \lean{GaussianFilter.continuous_intertwinerIntegrand,
        GaussianFilter.norm_intertwinerIntegrand,
        GaussianFilter.integrable_intertwinerIntegrand,
        GaussianFilter.integrable_intertwinerIntegrand_restrict}
    \leanok
    \uses{def:gaussian_filter_integrals}
    The function $t\mapsto A_t$ is continuous for arbitrary matrices
    $\mathbb H',\mathbb H,W$. If $\mathbb H'$ and $\mathbb H$ are
    Hermitian, then $\|A_t\|=\|W\|$ for every $t$.
    In that case $A_t$ is integrable against $\mu_h$ and its restriction
    to $[-T,T]$, for every $h\ge0$ and every $T\in\R$.
\end{theorem}
\begin{proof}\leanok
    \uses{def:gaussian_filter_integrals}
    Continuity follows from continuity of the matrix exponential and
    multiplication. For Hermitian generators, both exponential factors
    are unitary, so left and right multiplication preserve the operator
    norm. The constant $\|W\|$ is integrable under the probability
    measure $\mu_h$, and also under its restriction.
\end{proof}

\begin{theorem}[Density formulas and zero variance]
    \label{thm:gaussian_filter_density}
    \lean{GaussianFilter.gaussianIntertwiner_eq_integral_kernel,
        GaussianFilter.gaussianIntertwinerTruncated_eq_integral_kernel,
        GaussianFilter.integrable_kernel_intertwinerIntegrand,
        GaussianFilter.gaussianIntertwiner_zero_variance,
        GaussianFilter.gaussianIntertwiner_zero}
    \leanok
    \uses{def:gaussian_filter_integrals}
    For $h>0$ and arbitrary matrices $\mathbb H',\mathbb H,W$,
    \begin{align}
        M_h&=\int_{\R}g_h(t)e^{it\mathbb H'}We^{-it\mathbb H}\,dt,
        \label{eq:gaussian_filter_density}\\
        M_h^T&=\int_{[-T,T]}g_h(t)e^{it\mathbb H'}We^{-it\mathbb H}\,dt. \notag
    \end{align}
    If the two generators are Hermitian, these density integrands are
    integrable. At zero variance, $M_0(\mathbb H',\mathbb H;W)=W$.
    For every $h\ge0$, $M_h(\mathbb H',\mathbb H;0)=0$.
\end{theorem}
\begin{proof}\leanok
    \uses{thm:gaussian_density_integrals, thm:gaussian_integrand,
        def:gaussian_filter_integrals}
    Apply Theorem~\ref{thm:gaussian_density_integrals} to $A_t$ and its
    restriction; integrability follows from
    Theorem~\ref{thm:gaussian_integrand}. At zero variance, integration
    against $\delta_0$ evaluates $A_0=W$. For zero input, $A_t=0$.
\end{proof}

\begin{theorem}[Contraction and truncation error]
    \label{thm:gaussian_filter_contraction}
    \lean{GaussianFilter.norm_gaussianIntertwiner_le,
        GaussianFilter.norm_gaussianIntertwinerTruncated_le_mass,
        GaussianFilter.norm_gaussianIntertwinerTruncated_le,
        GaussianFilter.norm_gaussianIntertwiner_sub_truncated_le_mass,
        GaussianFilter.norm_gaussianIntertwiner_sub_truncated_le}
    \leanok
    \uses{def:gaussian_filter_integrals}
    Suppose $\mathbb H'$ and $\mathbb H$ are Hermitian and $h\ge0$.
    For every $T\in\R$,
    \begin{align}
        \|M_h\|&\le\|W\|, &
        \|M_h^T\|&\le\|W\|\mu_h([-T,T])\le\|W\|,
        \label{eq:gaussian_contractions}\\
        \|M_h-M_h^T\|&\le\mu_h([-T,T]^c)\|W\|. \notag
    \end{align}
    In particular, for $T\ge0$,
    \begin{align}
        \|M_h-M_h^T\|&\le\tau_h(T)\|W\|.
        \label{eq:gaussian_truncation_error}
    \end{align}
\end{theorem}
\begin{proof}\leanok
    \uses{thm:gaussian_integrand, thm:gaussian_tail}
    Integrating $\|A_t\|=\|W\|$ gives the full and restricted norm
    bounds. Since $A_t$ is integrable, the difference $M_h-M_h^T$ is
    its integral over $[-T,T]^c$. Bound its norm by the integral of
    $\|W\|$ and apply Theorem~\ref{thm:gaussian_tail}.
\end{proof}

\begin{theorem}[Adjoints exchange the generators]
    \label{thm:gaussian_filter_adjoint}
    \lean{GaussianFilter.gaussianIntertwiner_conjTranspose,
        GaussianFilter.gaussianIntertwinerTruncated_conjTranspose}
    \leanok
    \uses{def:gaussian_filter_integrals}
    For Hermitian $\mathbb H',\mathbb H$, every $h\ge0$ and $T\in\R$,
    \begin{align}
        M_h(\mathbb H',\mathbb H;W)^\dagger
            &=M_h(\mathbb H,\mathbb H';W^\dagger),
        \label{eq:gaussian_adjoint}\\
        M_h^T(\mathbb H',\mathbb H;W)^\dagger
            &=M_h^T(\mathbb H,\mathbb H';W^\dagger). \notag
    \end{align}
\end{theorem}
\begin{proof}\leanok
    \uses{def:gaussian_filter_integrals}
    Taking the adjoint commutes with the integral, and
    $A_t^\dagger=e^{it\mathbb H}W^\dagger e^{-it\mathbb H'}$.
    This identity holds for the full measure and for its restriction.
\end{proof}

\begin{theorem}[Matrix elements and the Gaussian multiplier]
    \label{thm:gaussian_filter_multiplier}
    \lean{GaussianFilter.dotProduct_intertwinerIntegrand_mulVec,
        GaussianFilter.dotProduct_gaussianIntertwiner_mulVec,
        GaussianFilter.inner_gaussianIntertwiner}
    \leanok
    \uses{def:gaussian_filter_integrals}
    Let $\mathbb H'v=Ev$ and $\mathbb Hw=E'w$, where $E,E'\in\R$.
    If $\mathbb H'$ is Hermitian, then
    \begin{align}
        \langle v,A_tw\rangle
            &=e^{it(E-E')}\langle v,Ww\rangle.
        \label{eq:gaussian_integrand_coefficient}
    \end{align}
    If both generators are Hermitian, then for every $h\ge0$,
    \begin{align}
        \langle v,M_hw\rangle
            &=e^{-h(E-E')^2/2}\langle v,Ww\rangle.
        \label{eq:gaussian_multiplier}
    \end{align}
    The vectors need not be normalized or real. In coordinates,
    $\langle v,M_hw\rangle=\sum_j\overline{v_j}(M_hw)_j$.
\end{theorem}
\begin{proof}\leanok
    \uses{thm:gaussian_integrand, thm:gaussian_characteristic}
    The right eigenvector contributes $e^{-itE'}$; the left covector
    contributes $e^{itE}$, because the first argument of the inner
    product is conjugate-linear. Their product proves
    (\ref{eq:gaussian_integrand_coefficient}). The continuous linear
    functional $A\mapsto\langle v,Aw\rangle$ commutes with the
    integrable matrix average. Apply
    (\ref{eq:gaussian_characteristic}) at $\omega=E-E'$ to obtain
    (\ref{eq:gaussian_multiplier}).
\end{proof}

\subsection{Spectral gaps and ground-vector estimates}

\begin{theorem}[Ground and excited eigenvectors]
    \label{thm:gaussian_gap_dichotomy}
    \lean{GaussianFilter.eigenvector_ground_or_gap}
    \leanok
    Let $K$ be Hermitian, $\|\Omega\|=1$, $K\Omega=\mathcal E\Omega$,
    and $\Delta>0$. Suppose
    \begin{align}
        K-\mathcal E I&\ge\Delta(I-\ketbra{\Omega}{\Omega}).
        \label{eq:gaussian_lower_gap}
    \end{align}
    For an orthonormal eigenbasis $Ku_k=\lambda_ku_k$, each $k$
    satisfies one of the following alternatives:
    \begin{align}
        \lambda_k&=\mathcal E, &
        u_k&=\langle\Omega,u_k\rangle\Omega;
        \label{eq:gaussian_eigenvector_cases}\\
        \lambda_k-\mathcal E&\ge\Delta, &
        \langle u_k,\Omega\rangle&=0. \notag
    \end{align}
    Thus $\Delta$ is a lower bound on the gap; equality with the least
    excitation energy is not required.
\end{theorem}
\begin{proof}\leanok
    When $\lambda_k=\mathcal E$, the quadratic-form inequality
    (\ref{eq:gaussian_lower_gap}) forces
    $\|u_k-\langle\Omega,u_k\rangle\Omega\|^2=0$.
    Otherwise Hermiticity and the two eigenvector equations imply
    $(\lambda_k-\mathcal E)\langle u_k,\Omega\rangle=0$.
    Hence $u_k\perp\Omega$, and evaluating
    (\ref{eq:gaussian_lower_gap}) on $u_k$ gives
    $\lambda_k-\mathcal E\ge\Delta$.
\end{proof}

\begin{theorem}[Dimension-free decay from spectral coefficients]
    \label{thm:gaussian_coefficient_decay}
    \lean{GaussianFilter.norm_sub_ground_le_of_coefficients}
    \leanok
    Under the hypotheses of Theorem~\ref{thm:gaussian_gap_dichotomy},
    let $h\ge0$ and suppose $x,y\in\mathcal H$ satisfy
    \begin{align}
        \langle u_k,y\rangle
            &=e^{-h(\lambda_k-\mathcal E)^2/2}\langle u_k,x\rangle
        \label{eq:gaussian_spectral_coefficients}
    \end{align}
    for every $k$. Then
    \begin{align}
        \|y-\langle\Omega,x\rangle\Omega\|
            &\le e^{-h\Delta^2/2}\|x\|.
        \label{eq:gaussian_coefficient_decay}
    \end{align}
\end{theorem}
\begin{proof}\leanok
    \uses{thm:gaussian_gap_dichotomy}
    Put $r=y-\langle\Omega,x\rangle\Omega$ and
    $q=e^{-h\Delta^2/2}$. The ground coefficients of $r$ vanish.
    Each excited coefficient obeys
    $|\langle u_k,r\rangle|\le q|\langle u_k,x\rangle|$, since
    $\lambda_k-\mathcal E\ge\Delta>0$ and $h\ge0$.
    Parseval's identity gives
    \begin{align}
        \|r\|^2
            &=\sum_k|\langle u_k,r\rangle|^2
            \le q^2\sum_k|\langle u_k,x\rangle|^2
            =q^2\|x\|^2.
        \label{eq:gaussian_parseval}
    \end{align}
    Taking square roots proves~(\ref{eq:gaussian_coefficient_decay}),
    with no factor depending on $\dim\mathcal H$.
\end{proof}

Let $\mathbb H',\mathbb H$ now be Hermitian, with unit vectors
$\Phi,\Psi$ and common eigenvalue $\mathcal E$:
\begin{align}
    \mathbb H'\Phi&=\mathcal E\Phi, &
    \mathbb H\Psi&=\mathcal E\Psi.
    \label{eq:gaussian_ground_vectors}
\end{align}
For any operator $W$, put $z=\langle\Phi,W\Psi\rangle\in\C$ and
$q_h=e^{-h\Delta^2/2}$.
% The five integral-ground declarations passed native production, strict
% production, consumers, and five guarded axiom checks at revision 2ae6e259.
% Their source bytes are unchanged in the render snapshot at 63a3390.

\begin{theorem}[Forward ground-vector estimate]
    \label{thm:gaussian_ground_estimate}
    \lean{GaussianFilter.norm_gaussianIntertwiner_ground_le}
    \leanok
    \uses{def:gaussian_filter_integrals}
    Assume~(\ref{eq:gaussian_ground_vectors}), $\Delta>0$, and
    $\mathbb H'-\mathcal E I\ge\Delta(I-\ketbra{\Phi}{\Phi})$.
    For every $h\ge0$,
    \begin{align}
        \|M_h\Psi-z\Phi\|&\le q_h\|W\|.
        \label{eq:gaussian_ground_estimate}
    \end{align}
    Only the target Hamiltonian $\mathbb H'$ needs a gap for this bound.
\end{theorem}
\begin{proof}\leanok
    \uses{thm:gaussian_filter_multiplier, thm:gaussian_coefficient_decay}
    In an orthonormal eigenbasis $\mathbb H'u_k=\lambda_ku_k$,
    (\ref{eq:gaussian_multiplier}) gives
    \begin{align}
        \langle u_k,M_h\Psi\rangle
            &=e^{-h(\lambda_k-\mathcal E)^2/2}
                \langle u_k,W\Psi\rangle.
        \label{eq:gaussian_ground_coefficients}
    \end{align}
    Apply Theorem~\ref{thm:gaussian_coefficient_decay} with
    $x=W\Psi$, $y=M_h\Psi$, and $\Omega=\Phi$.
    Since $\|W\Psi\|\le\|W\|\|\Psi\|=\|W\|$, this proves
    (\ref{eq:gaussian_ground_estimate}). Thus the spectral coefficient
    identity is obtained from the integral itself.
\end{proof}

\begin{theorem}[Two-sided ground-vector estimate]
    \label{thm:gaussian_two_sided}
    \lean{GaussianFilter.gaussianIntertwiner_two_sided_ground_estimate}
    \leanok
    \uses{def:gaussian_filter_integrals}
    Assume~(\ref{eq:gaussian_ground_vectors}), $\Delta>0$, and both
    lower-gap inequalities
    \begin{align}
        \mathbb H-\mathcal E I&\ge\Delta(I-\ketbra{\Psi}{\Psi}),
        \label{eq:gaussian_two_gaps}\\
        \mathbb H'-\mathcal E I&\ge\Delta(I-\ketbra{\Phi}{\Phi}). \notag
    \end{align}
    For every $h\ge0$,
    \begin{align}
        \|M_h\Psi-z\Phi\|&\le q_h\|W\|, &
        \|M_h^\dagger\Phi-\overline z\Psi\|&\le q_h\|W\|.
        \label{eq:gaussian_two_sided}
    \end{align}
\end{theorem}
\begin{proof}\leanok
    \uses{thm:gaussian_ground_estimate, thm:gaussian_filter_adjoint}
    Apply Theorem~\ref{thm:gaussian_ground_estimate} once to
    $(\mathbb H',\mathbb H,W,\Phi,\Psi)$ and once to
    $(\mathbb H,\mathbb H',W^\dagger,\Psi,\Phi)$.
    By~(\ref{eq:gaussian_adjoint}), the second filtered operator is
    $M_h^\dagger$. Its overlap is
    $\langle\Psi,W^\dagger\Phi\rangle
      =\overline{\langle\Phi,W\Psi\rangle}=\overline z$, and
    $\|W^\dagger\|=\|W\|$.
\end{proof}

\begin{theorem}[Real overlap and contraction]
    \label{thm:gaussian_real_overlap}
    \lean{GaussianFilter.gaussianIntertwiner_two_sided_of_real_overlap}
    \leanok
    \uses{def:gaussian_filter_integrals}
    Under~(\ref{eq:gaussian_ground_vectors}) and
    (\ref{eq:gaussian_two_gaps}), suppose $\Delta>0$, $\|W\|\le1$,
    and $\langle\Phi,W\Psi\rangle=z\in\R$.
    For every $h\ge0$,
    \begin{align}
        \|M_h\|&\le1, &
        \|M_h\Psi-z\Phi\|&\le q_h, &
        \|M_h^\dagger\Phi-z\Psi\|&\le q_h.
        \label{eq:gaussian_real_overlap}
    \end{align}
\end{theorem}
\begin{proof}\leanok
    \uses{thm:gaussian_two_sided, thm:gaussian_filter_contraction}
    Use Theorem~\ref{thm:gaussian_filter_contraction} and
    (\ref{eq:gaussian_two_sided}), then substitute
    $\overline z=z$ and $\|W\|\le1$.
\end{proof}

In the reset argument of~\cite{OpenAI2026PolynomialPEPS}, the overlap
$z$ is a purity and is therefore real and nonnegative. This is the
real-overlap hypothesis used in~(\ref{eq:gaussian_real_overlap});
Hermiticity of $W$ alone does not make $\langle\Phi,W\Psi\rangle$
real for two different vectors.

\begin{theorem}[Forward estimate after time truncation]
    \label{thm:gaussian_truncated_ground}
    \lean{GaussianFilter.norm_gaussianIntertwinerTruncated_ground_le}
    \leanok
    \uses{def:gaussian_filter_integrals}
    Under the hypotheses of Theorem~\ref{thm:gaussian_ground_estimate},
    for $h\ge0$ and $T\ge0$,
    \begin{align}
        \|M_h^T\Psi-z\Phi\|
            &\le\bigl(q_h+\tau_h(T)\bigr)\|W\|.
        \label{eq:gaussian_truncated_ground}
    \end{align}
\end{theorem}
\begin{proof}\leanok
    \uses{thm:gaussian_ground_estimate, thm:gaussian_filter_contraction}
    Split $M_h^T\Psi-z\Phi=(M_h\Psi-z\Phi)+(M_h^T-M_h)\Psi$.
    The triangle inequality, $\|\Psi\|=1$, and
    (\ref{eq:gaussian_truncation_error}) bound the two terms by
    $q_h\|W\|$ and $\tau_h(T)\|W\|$.
\end{proof}

\begin{theorem}[Two-sided estimate after time truncation]
    \label{thm:gaussian_truncated_two_sided}
    \lean{GaussianFilter.gaussianIntertwinerTruncated_two_sided_ground_estimate}
    \leanok
    \uses{def:gaussian_filter_integrals}
    Under~(\ref{eq:gaussian_ground_vectors}) and
    (\ref{eq:gaussian_two_gaps}), suppose $\Delta>0$, $h\ge0$, and $T\ge0$.
    With $\varepsilon=(q_h+\tau_h(T))\|W\|$,
    \begin{align}
        \|M_h^T\Psi-z\Phi\|&\le\varepsilon, &
        \|(M_h^T)^\dagger\Phi-\overline z\Psi\|&\le\varepsilon.
        \label{eq:gaussian_truncated_two_sided}
    \end{align}
    In particular, at positive variance the common bound is exactly
    \begin{align}
        \varepsilon&=\left(e^{-h\Delta^2/2}
            +2e^{-T^2/(2h)}\right)\|W\|.
        \label{eq:gaussian_truncated_error}
    \end{align}
\end{theorem}
\begin{proof}\leanok
    \uses{thm:gaussian_truncated_ground, thm:gaussian_filter_adjoint}
    Apply Theorem~\ref{thm:gaussian_truncated_ground} to the original
    and exchanged generators. The truncated identity in
    (\ref{eq:gaussian_adjoint}) converts the second result to the
    adjoint estimate, with overlap $\overline z$ and unchanged norm
    $\|W^\dagger\|=\|W\|$.
\end{proof}
